\section{Summary and Future Directions}
\label{sec:flex_summary}

This chapter presented \FLEX{}, an object-centric force-space framework for articulated object manipulation. \FLEX{} learns force policies in simulation without modeling a robot, estimates an object's joint structure through physical interaction, and maps the desired forces to robot commands through an operational space controller. The experiments demonstrated transfer to unseen articulated objects and across three robotic platforms in simulation without policy retraining, together with a real-world demonstration on a physical UR5 robot.

\FLEX{} instantiates the interaction abstraction framework developed in \cref{sec:interaction_abstractions}. Its generalization rests on exploiting structure inherent in the manipulation task. Objects that differ in appearance or purpose can nevertheless present the same underlying interaction: a drawer and a dishwasher rack both move along a prismatic joint, just as different doors can rotate about revolute joints. By representing this articulation structure and defining actions as forces on the object, \FLEX{} preserves what determines successful motion while abstracting away object identity and embodiment-specific control. In the terminology of \cref{ch:novelty_open_world}, this expands the range of variation that can be assimilated: the evaluated changes in object instance and robot embodiment are handled within the existing policy representation rather than through policy retraining. The result supports the dissertation's broader claim that an abstraction influences both what an agent must represent and what it can safely ignore.

Several limitations define the boundary of this result. \FLEX{} assumes quasi-static interaction and that a mechanism belongs to one of two supported classes: prismatic or revolute. Its real-world evaluation is also limited to a single robot--object demonstration. Dynamic manipulation, mechanisms with more complex articulation, and changing contact modes would require richer interaction representations and execution mappings. The framework further treats each articulated interaction as an individual skill. Integrating \FLEX{} policies as executors within the planning--learning framework of \cref{ch:rapid_learn} would allow these transferable skills to be composed into longer-horizon plans and would provide a mechanism for targeted accommodation when force-space execution does fail.

The next chapter examines whether the same interaction-abstraction principle extends beyond articulated mechanisms to rigid-body transport under unknown physical conditions.
